\chapter{Multi-Layer Perceptrons \& Backpropagation}\label{ch:mlp}
\begin{center}\small\textit{One line cannot solve XOR; stacking lines can approximate anything -- if we can train them.}\end{center}

\section{Intuition: Why Go Deeper?}
A single perceptron draws one line. XOR needs two lines combined: ``$(x_1\lor x_2)$ but not $(x_1\land x_2)$''. That is a hidden layer. Stack perceptrons: first layer draws several lines, second layer combines them into convex regions, third layer into arbitrary shapes. This is the \textbf{Multi-Layer Perceptron (MLP)} -- the simplest deep network.

The obstacle that froze the field for 15 years: how to assign blame to hidden weights when only the output error is visible? The answer is the \textbf{chain rule} applied systematically -- backpropagation.

\section{Formal Setup: Forward Pass}
An $L$-layer MLP maps $\mathbf{a}^{(0)}=\mathbf{x}\in\mathbb{R}^{d_0}$ to output $\hat{\mathbf{y}}=\mathbf{a}^{(L)}$ via
\begin{equation}
\mathbf{z}^{(l)} = W^{(l)}\mathbf{a}^{(l-1)}+\mathbf{b}^{(l)},\qquad \mathbf{a}^{(l)}=\sigma(\mathbf{z}^{(l)}),\qquad l=1,\dots,L
\end{equation}
where $W^{(l)}\in\mathbb{R}^{d_l\times d_{l-1}}$, $\sigma$ is a non-linear activation (sigmoid, tanh, ReLU). Without $\sigma$, the composition collapses to one affine map -- depth would be pointless.

Common activations:
\begin{equation}
\sigma_{\text{sig}}(z)=\frac{1}{1+e^{-z}},\quad \tanh(z)=\frac{e^{z}-e^{-z}}{e^{z}+e^{-z}},\quad \text{ReLU}(z)=\max(0,z)
\end{equation}

\begin{definition}[Universal Approximation, informal]
An MLP with one hidden layer and enough neurons can approximate any continuous function on a compact set arbitrarily well (Cybenko, 1989). Depth makes this efficient.
\end{definition}

\begin{center}
\begin{tikzpicture}[font=\footnotesize, >=Stealth, node distance=0.6cm]
  \tikzstyle{neuron}=[circle, draw, thick, minimum size=0.9cm, inner sep=0pt]
  % Input
  \node[neuron, fill=blue!15] (i1) at (0,1.5) {$x_1$};
  \node[neuron, fill=blue!15] (i2) at (0,0) {$x_2$};
  \node[neuron, fill=blue!15] (i3) at (0,-1.5) {$x_3$};
  % Hidden 1
  \node[neuron, fill=green!15] (h1) at (3,1.8) {$h_1$};
  \node[neuron, fill=green!15] (h2) at (3,0) {$h_2$};
  \node[neuron, fill=green!15] (h3) at (3,-1.8) {$h_3$};
  % Hidden 2
  \node[neuron, fill=orange!15] (g1) at (6,0.9) {$g_1$};
  \node[neuron, fill=orange!15] (g2) at (6,-0.9) {$g_2$};
  % Output
  \node[neuron, fill=violet!15] (o1) at (9,0) {$\hat y$};
  \foreach \i in {1,2,3} \foreach \h in {1,2,3} \draw[->, gray!60, thin] (i\i) -- (h\h);
  \foreach \h in {1,2,3} \foreach \g in {1,2} \draw[->, green!50!black, thin] (h\h) -- (g\g);
  \foreach \g in {1,2} \draw[->, orange!70!black, thick] (g\g) -- (o1);
  \node[above, font=\scriptsize, text=blue!60!black] at (0,2.3) {Input};
  \node[above, font=\scriptsize, text=green!50!black] at (3,2.6) {Hidden 1};
  \node[above, font=\scriptsize, text=orange!70!black] at (6,1.7) {Hidden 2};
  \node[above, font=\scriptsize, text=violet!70!black] at (9,0.8) {Output};
  \node[draw, fill=yellow!10, rounded corners, font=\tiny, align=center] at (3,-3.0) {$W^{(1)}\mathbf{x}+\mathbf{b}^{(1)}$\\{\scriptsize then }$\sigma$};
\end{tikzpicture}
\end{center}

\section{The Chain Rule and Gradients}
Training minimises a loss, e.g., MSE $\mathcal{L}=\tfrac12\|\mathbf{a}^{(L)}-\mathbf{y}\|^2$ or cross-entropy. We need $\partial\mathcal{L}/\partial W^{(l)}$ to step downhill.

For scalar composition $f(g(h(x)))$: $df/dx = f'(g)\,g'(h)\,h'(x)$. For vectors, Jacobians multiply. The cleverness of backprop is \emph{reusing} downstream gradients.

Define error signal $\boldsymbol{\delta}^{(l)}:=\partial\mathcal{L}/\partial\mathbf{z}^{(l)}$. Then:

\begin{align}
\boldsymbol{\delta}^{(L)} &= \nabla_{\mathbf{a}^{(L)}}\mathcal{L}\odot \sigma'(\mathbf{z}^{(L)})\\
\boldsymbol{\delta}^{(l)} &= \big(W^{(l+1)\top}\boldsymbol{\delta}^{(l+1)}\big)\odot \sigma'(\mathbf{z}^{(l)}),\quad l=L-1,\dots,1\\
\frac{\partial\mathcal{L}}{\partial W^{(l)}} &= \boldsymbol{\delta}^{(l)}\mathbf{a}^{(l-1)\top},\qquad \frac{\partial\mathcal{L}}{\partial\mathbf{b}^{(l)}}=\boldsymbol{\delta}^{(l)}
\end{align}
where $\odot$ is element-wise product. One forward pass caches $\mathbf{z}^{(l)},\mathbf{a}^{(l)}$; one backward pass propagates $\boldsymbol{\delta}$.

\begin{center}
\begin{tikzpicture}[font=\scriptsize, >=Stealth]
  \node[draw, fill=blue!12, rounded corners, minimum width=2.2cm, minimum height=0.9cm, align=center] (fwd) at (0,0) {Forward\\$\mathbf{a}^{(l-1)}\to\mathbf{z}^{(l)}\to\mathbf{a}^{(l)}$};
  \node[draw, fill=orange!12, rounded corners, minimum width=2.2cm, minimum height=0.9cm, align=center] (bwd) at (6,0) {Backward\\$\boldsymbol{\delta}^{(l+1)}\to\boldsymbol{\delta}^{(l)}$};
  \draw[->, thick, blue!60!black] (fwd) -- node[above] {cache $\mathbf{z}^{(l)}$} (bwd);
  \draw[->, thick, orange!70!black, bend left=30] (bwd) to node[below] {$\frac{\partial\mathcal{L}}{\partial W^{(l)}}$} (fwd);
  \node[draw, fill=green!10, rounded corners, font=\tiny, align=center] at (3,1.4) {Chain rule: re-use\\downstream gradient};
  \draw[->, dashed, violet!60!black] (3,1.1) -- (3,0.5);
\end{tikzpicture}
\end{center}

\section{Backpropagation Algorithm}
\begin{enumerate}[leftmargin=1.2em]
\item Forward: compute and store all $\mathbf{z}^{(l)},\mathbf{a}^{(l)}$.
\item Output error: $\boldsymbol{\delta}^{(L)} = (\mathbf{a}^{(L)}-\mathbf{y})\odot\sigma'(\mathbf{z}^{(L)})$ (for MSE).
\item Backward: for $l=L-1$ down to $1$, $\boldsymbol{\delta}^{(l)}=(W^{(l+1)\top}\boldsymbol{\delta}^{(l+1)})\odot\sigma'(\mathbf{z}^{(l)})$.
\item Gradients: $\partial\mathcal{L}/\partial W^{(l)}=\boldsymbol{\delta}^{(l)}\mathbf{a}^{(l-1)\top}$.
\item Update: $W^{(l)}\leftarrow W^{(l)}-\eta\,\partial\mathcal{L}/\partial W^{(l)}$.
\end{enumerate}

Complexity is $O(|W|)$ -- same as forward pass -- because each weight is touched once forward and once backward. Naive finite differences would cost $O(|W|^2)$.

\section{Worked Example: Tiny Network by Hand}
Network: $x\in\mathbb{R}$, one hidden neuron with sigmoid, linear output. $w_1=0.5,b_1=0$, $w_2=1.0,b_2=0$, input $x=1$, target $y=0.8$, loss $\tfrac12(\hat y-y)^2$, $\eta=0.5$. Forward: $z_1=0.5$, $a_1=\sigma(0.5)=0.622$, $\hat y=0.622$. Error: $\hat y-y=-0.178$. $\delta^{(2)}=-0.178$. $\partial\mathcal{L}/\partial w_2=\delta^{(2)}a_1=-0.111$. $\delta^{(1)}=\delta^{(2)}w_2\cdot\sigma'(z_1)=-0.178\times1\times0.235=-0.0418$. $\partial\mathcal{L}/\partial w_1=\delta^{(1)}x=-0.0418$. Updates: $w_2\leftarrow1.055$, $w_1\leftarrow0.521$. Loss drops from $0.0158$ to $0.0121$.

\begin{example}[Autograd in NumPy]
\begin{lstlisting}[language=Python]
import numpy as np
def sigmoid(z): return 1/(1+np.exp(-z))
def dsigmoid(z): s=sigmoid(z); return s*(1-s)

# 2-layer MLP: 2 -> 3 -> 1
np.random.seed(0)
W1 = np.random.randn(3,2)*0.5; b1=np.zeros(3)
W2 = np.random.randn(1,3)*0.5; b2=np.zeros(1)
X = np.array([[0,1],[1,0],[1,1],[0,0]],dtype=float).T  # (2,4)
Y = np.array([[1,1,0,0]],dtype=float)
lr=0.5
for epoch in range(500):
    # Forward
    Z1=W1@X+b1[:,None]; A1=np.maximum(0,Z1)  # ReLU
    Z2=W2@A1+b2[:,None]; A2=sigmoid(Z2)
    loss= -np.mean(Y*np.log(A2+1e-9)+(1-Y)*np.log(1-A2+1e-9))
    # Backward
    dZ2=A2-Y
    dW2=dZ2@A1.T/4; db2=dZ2.mean(axis=1)
    dA1=W2.T@dZ2; dZ1=dA1*(Z1>0)
    dW1=dZ1@X.T/4; db1=dZ1.mean(axis=1)
    W2-=lr*dW2; b2-=lr*db2; W1-=lr*dW1; b1-=lr*db1
print("Final loss",loss); print("Pred",A2.round(2))
\end{lstlisting}
\end{example}

\section{Computation Graph View}
\begin{center}
\begin{tikzpicture}[font=\scriptsize, >=Stealth]
  \node[draw, fill=blue!10, circle, minimum size=0.9cm] (x) at (0,0) {$x$};
  \node[draw, fill=green!12, rounded corners, minimum width=1.3cm, minimum height=0.8cm] (mul) at (2.2,0) {$\times w_1$};
  \node[draw, fill=orange!12, rounded corners, minimum width=1.0cm, minimum height=0.8cm] (sig) at (4.3,0) {$\sigma$};
  \node[draw, fill=violet!12, rounded corners, minimum width=1.3cm, minimum height=0.8cm] (mul2) at (6.4,0) {$\times w_2$};
  \node[draw, fill=red!12, circle, minimum size=0.9cm] (loss) at (8.5,0) {$\mathcal{L}$};
  \draw[->, thick, blue!60!black] (x) -- (mul);
  \draw[->, thick] (mul) -- (sig); \draw[->, thick] (sig) -- (mul2); \draw[->, thick] (mul2) -- (loss);
  % backward arrows
  \draw[->, thick, red!70!black, dashed] (loss) -- ++(0,-1.1) -- ++(-8.5,0) -- (x);
  \node[below, font=\tiny, text=red!70!black, align=center] at (4.3,-1.1) {backward: $\partial\mathcal{L}/\partial w$ flows via chain rule};
  \node[above, font=\tiny] at (1.1,0.3) {fwd};
  \node[above, font=\tiny] at (5.35,0.3) {fwd};
\end{tikzpicture}
\end{center}

\section{Initialization and Vanishing Gradients}
Backprop needs a starting point. Zero init fails (all neurons identical).
\textbf{Xavier/Glorot} sets $\operatorname{Var}[W]=2/(\text{fan}_{in}+\text{fan}_{out})$
for tanh; \textbf{He} sets $\operatorname{Var}[W]=2/\text{fan}_{in}$ for ReLU,
preserving variance forward and backward. Without this, deep nets see
exponentially vanishing or exploding activations.

Sigmoid/tanh saturate: $\sigma'(z)\approx0$ for $|z|>3$, so gradients die in
deep stacks. ReLU fixes this for $z>0$ ($\sigma'=1$) but can die ($z\le0$).
Leaky ReLU, GELU, and careful norm (BatchNorm, LayerNorm) mitigate.

\section{Computation Graph and Automatic Differentiation}
Modern frameworks build a directed acyclic graph: leaves are parameters and
inputs, internal nodes are ops ($+$, $\times$, $\sigma$). Forward evaluates and
caches; backward applies chain rule node-by-node (reverse-mode AD). Cost is
$\approx2\times$ forward, memory $O(\text{nodes})$. Gradient checking via finite
differences validates: $\partial\mathcal{L}/\partial w\approx[\mathcal{L}(w+\epsilon)-\mathcal{L}(w-\epsilon)]/(2\epsilon)$.

\section{Why Depth Beats Width}
A shallow net needs exponentially many hidden units to approximate a deep
compositional function (e.g., parity, or $f(x)=\sin(\sin(\dots\sin(x)))$).
Depth reuses features hierarchically: edges $\to$ motifs $\to$ objects. Formal
results (Eldan \& Shamir, 2016) exhibit functions where a 3-layer net needs
exponential width to match a 2-layer deep net. Optimization is harder with depth,
but residuals and normalization make it trainable -- foreshadowing ResNets
(Chapter~5).

\section{Regularization Preview}
Even with perfect optimization, MLPs overfit. Weight decay, dropout, and early
stopping (Chapter~4) inject bias toward simple functions. BatchNorm also
smooths the landscape, allowing larger learning rates -- a theme revisited in
Chapter~3.

\section{Common Pitfalls}
Forgetting to zero gradients (\texttt{optimizer.zero\_grad()} in PyTorch) causes
accumulation across batches -- loss explodes silently. Another trap: applying
softmax before cross-entropy (which already includes log-softmax) doubly squashes
and kills gradients. Use \texttt{CrossEntropyLoss} on logits. Finally, shuffling
matters: without shuffling, batches are correlated and BatchNorm statistics
drift.

\section{Further Reading}
Rumelhart, Hinton \& Williams (1986) introduced backprop for MLPs. Goodfellow
et al., \emph{Deep Learning}, Chapters 6--8, formalizes the graph view. For
hands-on AD, see Karpathy's \texttt{micrograd} (50 lines that build autograd
from scratch) -- stepping through it demystifies every $\boldsymbol{\delta}$.

\section{Chapter Summary}
MLPs compose perceptrons: $\mathbf{a}^{(l)}=\sigma(W^{(l)}\mathbf{a}^{(l-1)}+
\mathbf{b}^{(l)})$. Backprop reuses downstream gradients via
$\boldsymbol{\delta}^{(l)}=(W^{(l+1)\top}\boldsymbol{\delta}^{(l+1)})\odot
\sigma'(\mathbf{z}^{(l)})$, costing $O(|W|)$. Initialization and non-linearity
choice decide whether gradients flow. Depth yields exponential expressivity over
width. Next: how to walk downhill efficiently.

\smallskip
\noindent\textit{Key takeaways:} (1) forward caches $\mathbf{z},\mathbf{a}$;
(2) backward propagates $\boldsymbol{\delta}$; (3) chain rule gives
$\partial\mathcal{L}/\partial W^{(l)}=\boldsymbol{\delta}^{(l)}\mathbf{a}^{(l-1)\top}$;
(4) He/Xavier init prevents vanishing; (5) computation graphs unify all gradients.

\section{Exercises}
\begin{exercise} Derive $\boldsymbol{\delta}^{(l)}$ formula from first principles for MSE + sigmoid. Show each step of the Jacobian chain.\end{exercise}
\begin{exercise} Why does a deep stack of linear layers equal one linear layer? Prove $W_2(W_1\mathbf{x})=W\mathbf{x}$ and explain why non-linearity is essential.\end{exercise}
\begin{exercise} Compute one backprop step by hand for the XOR network: 2 inputs, 2 hidden ReLU, 1 output. Use given weights and input $(1,0)\to1$.\end{exercise}
\begin{exercise} Implement the NumPy MLP above but swap ReLU for tanh. Compare convergence speed and final loss over 500 epochs. Explain vanishing-gradient differences.\end{exercise}
\begin{exercise} Backprop needs $O(|W|)$ memory for cached $\mathbf{z}$. For a 10-layer network with 1000 neurons/layer, estimate floats stored. How does gradient checkpointing trade compute for memory?\end{exercise}
% End of Chapter
